\chapter{Overview}

This blueprint records the current formal development of \texttt{RudinAnalysis}.
The code follows the chapter order of Rudin's \emph{Principles of Mathematical Analysis},
starting with ordered fields and the construction of the real numbers, then moving to
metric topology and the theory of numerical sequences.

The project is intentionally elementary in style: many results are reproved directly from
the chosen axioms so that the formal development mirrors the book as closely as possible.

\begin{definition}
\label{def:project-scope}
The present formalization is organized into three mathematical blocks:
ordered algebraic structures, basic topology of metric spaces, and convergence theory for
sequences and subsequences.
\end{definition}

\chapter{The Real and Complex Number Systems}

\section{Ordered sets and completeness}

\begin{definition}
\label{def:ordered-bounds}
\lean{OrderedSets.is_upper_bound}
\lean{OrderedSets.is_lower_bound}
\lean{OrderedSets.is_bounded_above}
\lean{OrderedSets.is_bounded_below}
\lean{OrderedSets.is_bounded}
For $E \subseteq \alpha$ and $u,\ell \in \alpha$:
\[
u \text{ is an upper bound of } E \iff \forall x \in E,\ x \le u,
\]
\[
\ell \text{ is a lower bound of } E \iff \forall x \in E,\ \ell \le x.
\]
Accordingly,
\[
E \text{ is bounded above} \iff \exists u,\ \forall x \in E,\ x \le u,
\]
\[
E \text{ is bounded below} \iff \exists \ell,\ \forall x \in E,\ \ell \le x.
\]
\end{definition}

\begin{definition}
\label{def:sup-inf}
\lean{OrderedSets.is_supremum}
\lean{OrderedSets.is_infimum}
For $E \subseteq \alpha$ and $s,i \in \alpha$:
\[
s = \sup E \iff
\bigl(\forall x \in E,\ x \le s\bigr) \land
\bigl(\forall u,\ (\forall x \in E,\ x \le u) \to s \le u\bigr),
\]
\[
i = \inf E \iff
\bigl(\forall x \in E,\ i \le x\bigr) \land
\bigl(\forall \ell,\ (\forall x \in E,\ \ell \le x) \to \ell \le i\bigr).
\]
\end{definition}

\begin{definition}
\label{def:completeness-properties}
\lean{OrderedSets.have_least_upper_bound_property}
\lean{OrderedSets.have_greatest_lower_bound_property}
The completeness predicates are:
\[
\forall E \subseteq \alpha,\ E \neq \varnothing \land E \text{ bounded above }
\to \exists s,\ s = \sup E,
\]
\[
\forall E \subseteq \alpha,\ E \neq \varnothing \land E \text{ bounded below }
\to \exists i,\ i = \inf E.
\]
\end{definition}

\begin{theorem}
\label{thm:lub-glb-equivalent}
\lean{OrderedSets.bound_property_iff}
\uses{def:completeness-properties,def:sup-inf}
For a linear order,
\[
\text{LUB property} \iff \text{GLB property}.
\]
\end{theorem}

\begin{theorem}
\label{thm:inf-exists-from-complete}
\lean{OrderedSets.Completness.inf_exists}
\uses{def:completeness-properties,thm:lub-glb-equivalent}
Any type equipped with the project's completeness structure satisfies the
greatest-lower-bound property.
\end{theorem}

\begin{theorem}
\label{thm:sup-characterization}
\lean{OrderedSets.the_1_11_l}
\uses{def:sup-inf,thm:inf-exists-from-complete}
Every nonempty set that is bounded above admits a supremum.
\end{theorem}

\begin{theorem}
\label{thm:inf-characterization}
\lean{OrderedSets.the_1_11_r}
\uses{def:sup-inf,thm:inf-exists-from-complete}
Every nonempty set that is bounded below admits an infimum.
\end{theorem}

\section{Field axioms and ordered fields}

\begin{definition}
\label{def:field-axioms}
\lean{Fields.FieldAxioms}
The project introduces a custom structure of field axioms, close to Rudin's presentation,
in order to rederive elementary algebraic laws from first principles.
\end{definition}

\begin{theorem}
\label{thm:field-basic-cancellation}
\lean{Fields.FieldAxioms.add_left_cancel}
\lean{Fields.FieldAxioms.mul_left_cancel}
\uses{def:field-axioms}
Within any type satisfying these field axioms, one has additive and multiplicative
cancellation laws.
\end{theorem}

\begin{definition}
\label{def:ordered-field}
\lean{Fields.OrderedField}
An ordered field is a linear order together with the project's field axioms and the expected
compatibility of order with addition and multiplication by positive elements.
\end{definition}

\begin{theorem}
\label{thm:ordered-field-sign-laws}
\lean{Fields.OrderedField.neg_lt_zero_of_pos}
\lean{Fields.OrderedField.mul_self_pos}
\lean{Fields.OrderedField.inv_lt_inv_of_pos}
\uses{def:ordered-field,thm:field-basic-cancellation}
In an ordered field, negatives of positive elements are negative, squares of nonzero
elements are positive, and inversion reverses order on positive elements.
\end{theorem}

\section{Irrationality and Dedekind reals}

\begin{theorem}
\label{thm:sqrt-two-irrational}
\lean{Introduction.no_rational_sqrt_two}
The formal development proves that there is no rational number whose square is equal to~$2$.
\end{theorem}

\begin{definition}
\label{def:dedekind-real}
\lean{TheRealField.DedekindReal}
The real numbers are constructed as Dedekind cuts of rational numbers.
\end{definition}

\begin{theorem}
\label{thm:dedekind-complete}
\lean{TheRealField.instCompletnessDedekindReal}
\uses{def:dedekind-real}
The ordered set of Dedekind reals satisfies the completeness axiom used elsewhere in the
project.
\end{theorem}

\begin{theorem}
\label{thm:dedekind-additive-structure}
\lean{TheRealField.dedekindreal_add_comm}
\lean{TheRealField.dedekindreal_add_assoc}
\lean{TheRealField.dedekindreal_add_neg_cancel}
\uses{def:dedekind-real}
Addition on Dedekind reals is commutative and associative, and every element has an
additive inverse.
\end{theorem}

\begin{theorem}
\label{thm:dedekind-multiplicative-structure}
\lean{TheRealField.zero_lt_one}
\lean{TheRealField.mul_comm}
\lean{TheRealField.mul_assoc}
\lean{TheRealField.mul_inv_cancel}
\lean{TheRealField.mul_add}
\uses{def:dedekind-real,thm:dedekind-additive-structure}
The Dedekind reals carry the expected ordered-field operations: multiplication is
commutative and associative, distributes over addition, and every nonzero element admits a
multiplicative inverse.
\end{theorem}

\chapter{Basic Topology}

\section{Metric preliminaries}

\begin{definition}
\label{def:metric-balls}
\lean{MetricSpaces.closedBall}
\lean{MetricSpaces.openBall}
\lean{MetricSpaces.deletedBall}
For a metric space $(X,d)$, $x \in X$, and $r > 0$:
\[
B(x,r)=\{y\in X : d(x,y)<r\},
\]
\[
\overline{B}(x,r)=\{y\in X : d(x,y)\le r\},
\]
\[
B^\ast(x,r)=\{y\in X : 0<d(x,y)<r\}.
\]
\end{definition}

\begin{definition}
\label{def:bounded-limit-point}
\lean{MetricSpaces.IsBounded}
\lean{MetricSpaces.limitPoint}
\lean{MetricSpaces.sequenceBounded}
The basic notions are encoded as:
\[
E \text{ bounded } \iff \exists x\in X\ \exists R>0,\ E \subseteq B(x,R),
\]
\[
p \text{ is a limit point of } E
\iff \forall r>0,\ \bigl(B^\ast(p,r)\cap E\bigr)\neq\varnothing,
\]
\[
(x_n) \text{ bounded } \iff \exists x\in X\ \exists R>0,\ \forall n,\ x_n \in B(x,R).
\]
\end{definition}

\begin{theorem}
\label{thm:open-iff-ball-subset}
\lean{MetricSpaces.isOpen_iff_ball_subset}
\uses{def:metric-balls}
For $E \subseteq X$,
\[
E \text{ is open } \iff \forall x \in E,\ \exists r>0,\ B(x,r)\subseteq E.
\]
\end{theorem}

\begin{theorem}
\label{thm:open-ball-open}
\lean{MetricSpaces.isOpen_openBall}
\uses{def:metric-balls,thm:open-iff-ball-subset}
Every open ball is an open set.
\end{theorem}

\begin{theorem}
\label{thm:deleted-ball-open}
\lean{MetricSpaces.isOpen_deletedBall}
\uses{def:metric-balls,thm:open-iff-ball-subset}
Every deleted ball is open.
\end{theorem}

\begin{theorem}
\label{thm:limit-point-infinite-neighborhood}
\lean{MetricSpaces.limitPoint_ball_infinite}
\uses{def:metric-balls,def:bounded-limit-point}
If $p$ is a limit point of $E$, then
\[
\forall r>0,\ (B(p,r)\cap E)\text{ is infinite}.
\]
\end{theorem}

\section{Compactness}

\begin{definition}
\label{def:open-cover-data}
\lean{MetricSpaces.OpenCover.IsFinite}
\lean{MetricSpaces.SubCover.IsFinite}
\lean{MetricSpaces.SubCover.toOpenCover}
The project packages the finite-subcover language needed to compare Rudin's compactness with
the compactness API from \texttt{mathlib}.
\end{definition}

\begin{theorem}
\label{thm:compact-finite-subcover}
\lean{MetricSpaces.compact_iff_finite_subcover}
\lean{MetricSpaces.finite_subcover_of_compact}
\uses{def:open-cover-data}
For $K \subseteq X$,
\[
K \text{ compact }
\iff
\forall \{U_i\}_{i\in I},\ 
K \subseteq \bigcup_{i\in I} U_i \land (\forall i,\ U_i \text{ open})
\to \exists i_1,\dots,i_n,\ K \subseteq U_{i_1}\cup\cdots\cup U_{i_n}.
\]
\end{theorem}

\begin{theorem}
\label{thm:compact-induced-closed}
\lean{MetricSpaces.isCompact_induced_iff}
\lean{MetricSpaces.isClosed_of_isCompact}
\uses{thm:compact-finite-subcover}
Compactness behaves well under induced topologies, and compact subsets of metric spaces are
closed.
\end{theorem}

\begin{theorem}
\label{thm:finite-intersection-property}
\lean{MetricSpaces.nonempty_iInter_of_finite_intersections}
\uses{thm:compact-finite-subcover}
For a compact set, any family of closed subsets with the finite intersection property has
nonempty total intersection.
\end{theorem}

\begin{theorem}
\label{thm:bolzano-weierstrass}
\lean{MetricSpaces.bolzano_weierstrass}
\uses{def:bounded-limit-point,thm:compact-finite-subcover,thm:finite-intersection-property}
Every infinite subset of a compact metric space has a limit point in that compact set.
\end{theorem}

\chapter{Numerical Sequences and Series}

\section{Convergent sequences}

\begin{definition}
\label{def:sequence-convergence}
\lean{ConvergentSequences.convergesTo}
\lean{ConvergentSequences.convergent}
\lean{ConvergentSequences.Convergent}
\lean{ConvergentSequences.limit}
For a sequence $(x_n)$ in a metric space and $p \in X$,
\[
x_n \to p
\iff
\forall \varepsilon > 0,\ \exists N,\ \forall n \ge N,\ d(x_n,p)<\varepsilon.
\]
The blueprint also uses the bundled predicates ``$(x_n)$ is convergent'' and
``$\lim x_n$''.
\end{definition}

\begin{theorem}
\label{thm:convergence-eventually-ball}
\lean{ConvergentSequences.convergesTo_iff_eventually_in_ball}
\uses{def:sequence-convergence,def:metric-balls}
For a sequence $(x_n)$ and $p \in X$,
\[
x_n \to p
\iff
\forall r>0,\ \exists N,\ \forall n \ge N,\ x_n \in B(p,r).
\]
\end{theorem}

\begin{theorem}
\label{thm:limit-unique}
\lean{ConvergentSequences.limit_unique}
\uses{thm:convergence-eventually-ball}
Limits of sequences in metric spaces are unique.
\end{theorem}

\begin{theorem}
\label{thm:convergent-bounded}
\lean{ConvergentSequences.convergent_isBounded}
\uses{def:sequence-convergence,def:bounded-limit-point}
\[
x_n \to p \quad \Longrightarrow \quad (x_n)\text{ is bounded}.
\]
\end{theorem}

\begin{theorem}
\label{thm:sequence-from-limit-point}
\lean{ConvergentSequences.sequence_of_limitPoint}
\uses{def:bounded-limit-point,def:sequence-convergence}
If $p$ is a limit point of $E$, then there exists a sequence $(x_n)$ such that
\[
\forall n,\ x_n \in E,\qquad x_n \ne p,\qquad x_n \to p.
\]
\end{theorem}

\begin{theorem}
\label{thm:algebra-of-limits}
\lean{ConvergentSequences.cnv_add}
\lean{ConvergentSequences.cnv_const_mul}
\lean{ConvergentSequences.cnv_const_add}
\lean{ConvergentSequences.cnv_mul}
\uses{def:sequence-convergence,thm:limit-unique}
For complex-valued sequences,
\[
x_n \to x,\ y_n \to y \Longrightarrow x_n+y_n \to x+y,\quad x_ny_n \to xy,
\]
\[
x_n \to x,\ c \in \mathbb{C} \Longrightarrow cx_n \to cx,\quad x_n+c \to x+c.
\]
\end{theorem}

\section{Subsequences and compactness}

\begin{definition}
\label{def:subsequence}
\lean{Subsequences.SubSeq.val}
A subsequence of $(x_n)$ is given by a strictly increasing map
\[
\phi:\mathbb{N}\to\mathbb{N},\qquad \phi(m)<\phi(n)\ \text{ whenever }m<n,
\]
and the induced sequence is $(x_{\phi(n)})$.
\end{definition}

\begin{theorem}
\label{thm:subsequence-of-convergent}
\lean{Subsequences.subseq_converge_iff_converge}
\uses{def:subsequence,def:sequence-convergence}
If $x_n \to p$ and $\phi$ is strictly increasing, then
\[
x_{\phi(n)} \to p.
\]
\end{theorem}

\begin{theorem}
\label{thm:compact-has-convergent-subsequence}
\lean{Subsequences.compactSpace_has_convergent_subsequence}
\lean{Subsequences.bounded_sequence_in_euclidean_has_convergent_subsequence}
\uses{thm:bolzano-weierstrass,thm:sequence-from-limit-point}
If $(x_n)$ is a sequence in a compact metric space, then there exist a strictly increasing
$\phi:\mathbb{N}\to\mathbb{N}$ and $p$ such that
\[
x_{\phi(n)} \to p.
\]
In particular, every bounded sequence in $\mathbb{R}^m$ has a convergent subsequence.
\end{theorem}

\begin{definition}
\label{def:subsequential-limits}
\lean{Subsequences.subsequentialLimits}
The set of subsequential limits of $(x_n)$ is
\[
\operatorname{SubLim}(x_n)
=
\left\{
p \in X : \exists \phi:\mathbb{N}\to\mathbb{N}\text{ strictly increasing, }x_{\phi(n)}\to p
\right\}.
\]
\end{definition}

\begin{theorem}
\label{thm:subsequential-limits-closed}
\lean{Subsequences.subsequentialLimits_isClosed}
\uses{def:subsequential-limits,thm:compact-has-convergent-subsequence}
The set of subsequential limits of a sequence is closed.
\end{theorem}

\section{Cauchy sequences and completeness}

\begin{definition}
\label{def:cauchy-tail-complete}
\lean{Subsequences.cauchySequence}
\lean{Subsequences.tailSet}
\lean{Subsequences.completeMetricSpace}
For a sequence $(x_n)$ in a metric space:
\[
(x_n)\text{ is Cauchy }
\iff
\forall \varepsilon>0,\ \exists N,\ \forall m,n\ge N,\ d(x_m,x_n)<\varepsilon.
\]
Its $N$-th tail set is
\[
T_N=\{x_n : n\ge N\}.
\]
Completeness is expressed by
\[
X \text{ complete } \iff \text{every Cauchy sequence in }X\text{ converges}.
\]
\end{definition}

\begin{theorem}
\label{thm:cauchy-diameter-criterion}
\lean{Subsequences.cauchySequence_iff_diameter_tails_vanish}
\lean{Subsequences.diameter_closure_eq}
\uses{def:cauchy-tail-complete}
\[
(x_n)\text{ is Cauchy } \iff \operatorname{diam}(T_N)\to 0.
\]
Moreover,
\[
\operatorname{diam}(\overline{E})=\operatorname{diam}(E).
\]
\end{theorem}

\begin{theorem}
\label{thm:nested-compact-intersection}
\lean{Subsequences.singleton_intersection_of_nested_compact}
\uses{thm:finite-intersection-property,thm:cauchy-diameter-criterion}
If $(K_n)$ is a decreasing sequence of nonempty compact sets with
\[
K_{n+1}\subseteq K_n
\qquad\text{and}\qquad
\operatorname{diam}(K_n)\to 0,
\]
then
\[
\bigcap_{n=0}^\infty K_n=\{p\}
\]
for some point $p$.
\end{theorem}

\begin{theorem}
\label{thm:cauchy-convergence-results}
\lean{Subsequences.convergent_implies_cauchy}
\lean{Subsequences.cauchySequence_converges_in_compactSpace}
\lean{Subsequences.cauchySequence_converges_in_euclidean}
\uses{def:cauchy-tail-complete,thm:compact-has-convergent-subsequence,thm:limit-unique}
\[
x_n \to p \Longrightarrow (x_n)\text{ is Cauchy}.
\]
Also,
\[
X \text{ compact } \Longrightarrow \text{every Cauchy sequence in }X\text{ converges},
\]
and in particular every Cauchy sequence in $\mathbb{R}^m$ converges.
\end{theorem}

\begin{theorem}
\label{thm:completeness-results}
\lean{Subsequences.compactSpace_complete}
\lean{Subsequences.euclideanSpace_complete}
\lean{Subsequences.closed_subset_of_complete_is_complete}
\uses{def:cauchy-tail-complete,thm:cauchy-convergence-results}
\[
X \text{ compact } \Longrightarrow X \text{ complete},
\]
\[
\mathbb{R}^m \text{ is complete},
\]
\[
X \text{ complete and } E \subseteq X \text{ closed } \Longrightarrow E \text{ complete}.
\]
\end{theorem}

\section{Monotone sequences}

\begin{definition}
\label{def:monotone-sequences}
\lean{Subsequences.monotoneIncreasing}
\lean{Subsequences.monotoneDecreasing}
\lean{Subsequences.monotoneSequence}
For a real sequence $(s_n)$:
\[
\text{increasing} \iff \forall m\le n,\ s_m\le s_n,
\]
\[
\text{decreasing} \iff \forall m\le n,\ s_n\le s_m,
\]
\[
\text{monotone} \iff \text{increasing or decreasing}.
\]
\end{definition}

\begin{theorem}
\label{thm:monotone-convergent-iff-bounded}
\lean{Subsequences.monotoneSequence_convergent_iff_bounded}
\uses{def:monotone-sequences,thm:completeness-results}
For a monotone real sequence $(s_n)$,
\[
(s_n)\text{ converges } \iff (s_n)\text{ is bounded}.
\]
\end{theorem}
